\documentclass{beamer}
\usepackage[frenchb]{babel}
\usepackage[utf8]{inputenc}

\usepackage{tikz}

\input{gen/pygments-style.tex}
%\usepackage{xcolor}
\definecolor{MellonGreen}      {HTML}{D1F2A5}
\definecolor{MellonYellow}     {HTML}{EFFAB4}
\definecolor{MellonLightOrange}{HTML}{FFC48C}
\definecolor{MellonDarkOrange} {HTML}{FF9F80}
\definecolor{MellonPink}       {HTML}{F56991}

\usepackage{etoolbox}
\usepackage{xspace}
\usepackage{../switch}
\usepackage{../mathpartir}

\newcommand{\irule}[3]{ \inferrule*[right=(#1)]{#2}{#3} }

% Fix '::' -> ': :' when using babel
\newenvironment{grammarf}{\shorthandoff{:}\begin{grammar}}{\end{grammar}\shorthandon{:}}

\newcommand{\gramisa}[0]{::=~}
\newcommand{\gramor}[0]{|\hspace{3mm}}

\newcommand{\gramlr}[2]{%
  \begin{minipage}{0.2\textwidth}\raggedleft #1\end{minipage}
  \begin{minipage}{0.8\textwidth}#2\end{minipage}
}

% Types
\newcommand{\tInt}{\textsc{Int}\xspace}
\newcommand{\tFloat}{\textsc{Float}\xspace}
\newcommand{\tVoid}{\textsc{Void}\xspace}

\newcommand{\tFV}[1]{\textrm{FreeVars}(#1)}

\def\opxor{\mathbin{\char`\^}}

% Constant
\newcommand{\cNil}{\textsc{Nil}}

% Program & language names
\newcommand{\ctonewspeak}{\textsc{c2newspeak}}
\newcommand{\cpp}{\textsc{cpp}}
\newcommand{\newspeak}{\textsc{Newspeak}}

% Newspeak
\newcommand{\npkDecl}[2]{\textsc{Decl}~(#1) \{ #2 \}}
\newcommand{\npkIf}[3]{\textsc{If}(#1) \{ #2 \} \textsc{Else} \{ #3 \}}
\newcommand{\npkDoWith}[2]{\textsc{Do}\{ #1 \}\textsc{With}~#2:}
\newcommand{\npkWhile}[1]{\textsc{While}(1)\{#1\}}
\newcommand{\npkGoto}[1]{\textsc{Goto}~#1}
\newcommand{\npkReturn}[1]{\textsc{Return}~#1}

% Continuations
\def\kPass{\textsc{Kpass}}
\def\kSeq#1#2{\textsc{Kseq}(#1, #2)}
\def\kBlock#1#2{\textsc{Kblock}(#1, #2)}
\def\kGoto#1{\textsc{Kgoto}(#1)}
\def\kCall#1#2#3#4{\textsc{Kcall}(#1, #2, #3, #4)}
\def\kRet#1{\textsc{Kret}(#1)}

\newcommand{\insertcode}[1]{\input{gen/#1.pyg.tex}}

% \chapter* + add to toc
%\newcommand{\chapterstar}[1]{\chapter*{#1} \addcontentsline{toc}{chapter}{#1}}

% Register names
\def\defreg#1{\expandafter\def\csname#1\endcsname{\textsc{#1}\xspace}}

\defreg{eax}
\defreg{ebx}
\defreg{ecx}
\defreg{edx}
\defreg{esi}
\defreg{edi}
\defreg{ebp}
\defreg{esp}
\defreg{eip}
\defreg{eflags}

\def\crtrois{\textsc{CR3}}

\newcommand{\ring}[1]{\emph{ring}~#1}

\newcommand{\asminstr}[1]{\texttt{#1}}

\newcommand{\wip}[0]{%
\begin{center}
    \includegraphics[width=5cm]{gen/travaux.pdf}
\end{center}
}

\newcommand{\listcons}[0]{%
\texttt{(::)}%
}

\newcommand{\hexa}[1]{%
\texttt{0x#1}%
}

% Debug tikz bounding box
\def\drawbb{\draw[thick] (current bounding box.south west) rectangle (current bounding box.north east);}

% comes from the memoir manual {{{

\newcommand*{\addtotextheight}[1]{%
\twocolumn
\addtolength{\textheight}{#1}%
\onecolumn}

%% for main body, bottom of text at 1in, footer below
%% top of header at 1in, first text line double spaced
%% below base of header
\newlength{\linespace}
\setlength{\linespace}{\baselineskip} % current equivalent of \onelineskip
\setlength{\headheight}{\onelineskip}
\setlength{\headsep}{\linespace}
\addtolength{\headsep}{-\topskip}
\setlength{\uppermargin}{1in}
\addtolength{\uppermargin}{\headheight}
\addtolength{\uppermargin}{\headsep}

%% and for the bottom
\setlength{\lowermargin}{1in}
\setlength{\textheight}{\paperheight}

\addtolength{\textheight}{-\uppermargin}
\addtolength{\textheight}{-\lowermargin}
%% footnote settings
\setlength{\footskip}{\onelineskip}
\setlength{\footnotesep}{\onelineskip}


%% the fiddle lengths (..ta.. for title/approval page, others for prelims)
\newlength{\toptafiddle} \setlength{\toptafiddle}{2\linespace}
\newlength{\bottafiddle} \setlength{\bottafiddle}{0pt}
\newlength{\topfiddle}
\setlength{\topfiddle}{\toptafiddle}
\newlength{\botfiddle}
\setlength{\botfiddle}{\onelineskip}

%%% general macro for Abstract, etc., headings
\newcommand*{\pretoctitle}[1]{{\clearpage\centering \vspace*{-\topfiddle}#1\par}}

%%% Start the ACKNOWLEDGEMENTS
\newcommand{\asuacknowledgements}{\pretoctitle{ACKNOWLEDGEMENTS}}

%%% Start the ABSTRACT
\newcommand{\asuabstract}{ \addtotextheight{-\botfiddle} \pretoctitle{ABSTRACT}}

\addtodef{\mainmatter}{\addtotextheight{\botfiddle}{}}

%% make it easy to center any dedication
\newcommand{\asudedication}[1]{%
{\clearpage\mbox{}\vfill\centering #1 \par\vfill\clearpage}}

% }}} /memoir

\def\disprule#1{
  \switch{#1}
    {Cst-Int}{ \irule{Cst-Int}
                 { }
                 { Γ ⊢ i : \tInt}
    }
    {Cst-Float}{ \irule{Cst-Float}
                   { }
                   { Γ ⊢ d : \tFloat}
    }
    {Cst-Null}{ \irule{Cst-Null}
                  { }
                  { Γ ⊢ \eNull : t*}
    }
    {Cst-Unit}{ \irule{Cst-Unit}
                  { }
                  { Γ ⊢ \eUnit : \tUnit }
    }
    {Lv-Var}{ \irule{Lv-Var}
                { x : t ∈ Γ }
                { Γ ⊢ x : t }
    }
    {Lv-Deref}{ \irule{Lv-Deref}
                  { Γ ⊢ lv : t* }
                  { Γ ⊢ *lv : t }
    }
    {Lv-Index}{ \irule{Lv-Index}
                  { Γ ⊢ e : \tInt \\
                    Γ ⊢ lv : t[]
                  }
                  { Γ ⊢ lv[e] : t }
    }
    {Lv-Field}{ \irule{Lv-Field}
                  { (s, l, t_l) ∈ S \\
                    Γ ⊢ lv : \tstruct{s}
                  }
                  { Γ ⊢ lv.l : t_l }
    }
    {Op-Int}{ \irule{Op-Int}
                { \opbin ∈ \{+,-,\times,/,\&,|,\opxor,\&\&,||,\ll,\gg \} \\
                  Γ ⊢ e_1 : \tInt \\
                  Γ ⊢ e_2 : \tInt
                }{
                  Γ ⊢ e_1~\opbin~e_2 : \tInt
                }
    }
    {Op-Float}{ \irule{Op-Float}
                { \opbin ∈ \{+.,-.,\times.,/.\} \\
                  Γ ⊢ e_1 : \tFloat \\
                  Γ ⊢ e_2 : \tFloat
                }{
                  Γ ⊢ e_1~\opbin~e_2 : \tFloat
                }
    }
    {Eq-Num}{ \irule{Eq-Num}
                {t ∈ \{ \tInt, \tFloat \}}
                {\textsc{Eq}(t)}
    }
    {Eq-Ptr}{ \irule{Eq-Ptr}
                {\textsc{Eq}(t)}
                {\textsc{Eq}(t*)}
    }
    {Eq-Array}{ \irule{Eq-Array}
                  {\textsc{Eq}(t)}
                  {\textsc{Eq}(t[])}
    }
    {Op-Eq}{ \irule{Op-Eq}
             { \opbin ∈ \{=,≠\} \\
               Γ ⊢ e_1 : t \\
               Γ ⊢ e_2 : t \\
               \textsc{Eq}(t)
             }{
               Γ ⊢ e_1~\opbin~e_2 : \tInt
             }
    }
    {Op-Comparable}{ \irule{Op-Comparable}
                       { \opbin ∈ \{=,≠,≤,≥,<,>\} \\
                         Γ ⊢ e_1 : t \\
                         Γ ⊢ e_2 : t \\
                         t ∈ \{ \tInt, \tFloat \}
                       }{
                         Γ ⊢ e_1~\opbin~e_2 : \tInt
                       }
    }
    {Unop-Plus-Int}{ \irule{Unop-Plus-Int}
                       { Γ ⊢ e : \tInt }
                       { Γ ⊢ + e : \tInt }
    }
    {Unop-Plus-Float}{ \irule{Unop-Plus-Float}
                         { Γ ⊢ e : \tFloat }
                         { Γ ⊢ +. e : \tFloat }
    }
    {Unop-Minus-Int}{ \irule{Unop-Minus-Int}
                        { Γ ⊢ e : \tInt }
                        { Γ ⊢ - e : \tInt }
    }
    {Unop-Minus-Float}{ \irule{Unop-Minus-Float}
                          { Γ ⊢ e : \tFloat }
                          { Γ ⊢ -. e : \tFloat }
    }
    {Unop-Not}{ \irule{Unop-Not}
                { \opun ∈ \{ \sim, ! \} \\
                  Γ ⊢ e : \tInt
                }{
                  Γ ⊢ \opun~e : \tInt
                }
              }
    {Ptr-Arith}{ \irule{Ptr-Arith}
                   { \opbin ∈ \{ +_p, -_p \} \\
                     Γ ⊢ e_1 : t* \\
                     Γ ⊢ e_2 : \tInt \\
                   }
                   { Γ ⊢ e_1~\opbin~e_2 : t* }
               }
    {Addr}{ \irule{Addr}
              { Γ ⊢ lv : t }
              { Γ ⊢ \&lv : t* }
          }
    {Set}{ \irule{Set}
             { Γ ⊢ lv : t \\
               Γ ⊢ e : t
             }
             { Γ ⊢ lv ← e : t }
         }
    {Array}{ \irule{Array}
               { ∀ i ∈ [1;n] , Γ ⊢ e_i : t }
               { Γ ⊢ \{ e_1 ;
                       … ; e_n \}
                       : t[]
               }
           }
    {Struct}{ \irule{Struct}
                { ∀ i ∈ [1;n], Γ ⊢ e_i : t_i \\
                  ∀ i ∈ [1;n], (s, l_i, t_i) ∈ S
                }
                { Γ ⊢ \{ l_1 : e_1 ;
                     … ; l_n : e_n \}
                    : \tstruct{s}
                }
            }
    {Call}{ \irule{Call}
              { Γ ⊢ e : ( t_1 ,
                         … , t_n ) → t \\
                ∀ i ∈ [1;n], Γ ⊢ e_i : t_i
              }
              { Γ ⊢ e ( e_1
                         , …
                         , e_n
                         )
                      : t
              }
          }
    {Pass}{ \irule{Pass}
              { }
              {Γ ⊢ \iPass}
          }
    {Seq}{ \irule{Seq}
             { Γ ⊢ i_1 \\
               Γ ⊢ i_2
             }{
               Γ ⊢ i_1;i_2
             }
         }
    {Exp}{ \irule{Exp}
             { Γ ⊢ e : t }
             { Γ ⊢ e }
         }
    {If}{ \irule{If}
            { Γ ⊢ e : \tInt \\
              Γ ⊢ i_1 \\
              Γ ⊢ i_2
            }
            { Γ ⊢ \iIf{e}{i_1}{i_2} }
        }
    {While}{ \irule{While}
               { Γ ⊢ e : \tInt \\
                 Γ ⊢ i
               }
               { Γ ⊢ \iWhile{e}{i} }
           }
    {Return}{ \irule{Return}
                { Γ ⊢ \vRet ← e }
                { Γ ⊢ \iReturn{e} }
            }
    {Fun}{ \irule{Fun}
             { Γ' = (Γ - \vRet), a_1 : t_1, …, a_n : t_n,
                    l_1 : t'_1, …, l_n : t'_p,
                    \vRet : t \\
               ∀ i ∈ [1;p], Γ ⊢ e_i : t'_i \\
               Γ' ⊢ i
             }
             { Γ ⊢ \mathrm{fun} (a_1, …, a_n)
                   ((l_1, e_1), …, (l_p, e_p)) \{i\}
                    : (t_1, …, t_n) → t
             }
         }
    {Ph-Exp}{ \irule{Ph-Exp}
                { Γ ⊢ e : t }
                { \typh{Γ}{e}{Γ} }
            }
    {Ph-Var}{ \irule{Ph-Var}
                { Γ ⊢ e : t \\
                  Γ' = (x, t), Γ
                }
                { \typh{Γ}{p}{Γ'} }
            }
    {Ph-Struct}{ \irule{Ph-Struct}
                   { S' =
                     (x_1, t_1, s),
                     …
                     (x_n, t_n, s)
                   }
                   { \typh{Γ}{ \tstruct{s}
                                  \{ x_1 : t_1
                                   ; …
                                   ;  x_n : t_n
                                  \}
                              }{Γ}
                   }
               } % TODO
    {END}
}

\def\semrule#1{
  \switch{#1}
    {Ctx}{ \irule{Ctx}
             { \mm{m}{e}{m'}{e'} }
             { \mm{m}{\ctxSub{C}{e}}{m'}{\ctxSub{C}{e'}} }
    }
    {Ctx-Lv}{ \irule{Ctx-Lv}
             { \mmlv{m}{lv}{m'}{lv'} }
             { \mmlv{m}{\ctxSubL{C_L}{lv}}{m'}{\ctxSubL{C_L}{lv'}} }
    }
    {Ctx-Instr}{ \irule{Ctx-Instr}
             { \mmlv{m}{i}{m'}{i'} }
             { \mmlv{m}{\ctxSubI{C_I}{i}}{m'}{\ctxSubI{C_I}{i'}} }
    }
    {Phi-Var}{ \irule{Phi-Var}
                 { a = \mathrm{Lookup}(x, m) }
                 {
                   \mmlv{m}{x}{m}{a}
                 }
    }
    {Phi-Deref}{ \irule{Phi-Deref}
                   { }
                   { \mms{*φ}{\widehat{*}φ} }
    }
    {Phi-Struct}{ \irule{Phi-Struct}
                    {  }
                    { \mms{lv.l}{lv\widehat{.}l} }
    }
    {Phi-Array}{ \irule{Phi-Array}
                   { }
                   { \mms{φ[n]}{φ\widehat{[n]}} }
    }
    {Exp-Cst}{ \irule{Exp-Cst}
                 { }
                 {\mms{c}{\widehat{c}}}
             }
    {Exp-Fun}{ \irule{Exp-Fun}
                 { }
                 {\mms{f}{\widehat{f}}}
             }
    {Exp-Lv}{ \irule{Exp-Lv}
              { }
              {\mm{m}{φ}{m}{m[φ]_Φ}}
    }
    {Exp-UnOp}{ \irule{Exp-UnOp}
                { }
                {\mms{\opun~v}{\widehat{\opun}~v}}
    }
    {Exp-BinOp}{ \irule{Exp-BinOp}
                   { \opbin ∉ \{/, \%\} }
                   {\mms{v_1~\opbin~v_2}{v_1~\widehat{\opbin}~v_2}}
    }
    {Exp-Div}{ \irule{Exp-Div}
                 { \opbin ∈ \{/, \%\} \\
                   v_2 ≠ \widehat{0}
                 }
                 {\mms{v_1~\opbin~v_2}{v_1~\widehat{\opbin}~v_2}}
    }
    {Exp-Div-Zero}{ \irule{Exp-Div-Zero}
                 { \opbin ∈ \{/, \%\}
                 }
                 {\msi{m}{v_1~\opbin~0} → Ω_{div}}
    }
    {Exp-AddrOf}{ \irule{Exp-AddrOf}
                    { a = \mathrm{Lookup} (x, m)
                    }
                    {\mm{m}{\&x}{m}{a}}
    }
    {Exp-Set}{ \irule{Exp-Set}
                 { }
                 {\mm{m}{φ ← v}{m[φ ← v]}{v}}
    }
    {Exp-Struct}{ \irule{Exp-Struct}
                   { }
                   { \mms{ \{ l_1 : v_1 ; … ; l_n : v_n \} }
                         { \widehat{\{ l_1 : v_1 ; … ; l_n : v_n \}} }
                   }
    }
    {Exp-Array}{ \irule{Exp-Array}
                 { }
                 {
                   \mms{ [ v_1, … , v_n ] }
                       { \widehat{[ v_1, … , v_n ]} }
                 }
    }
    {Exp-Call}{ \irule{Exp-Call}
                  { f =
                    \mathrm{fun} (a_1, …, a_n)
                                 ((l'_1, e'_1), …, (l'_p, e'_p))
                                 \{ i \} \\
                    m_1 = \mathrm{Push}(m_0, ((a_1, v_1), …, (a_n, v_n))) \\
                    \evstar{m_1}{e'_1 \\ \vdots \\ e'_p}
                           {m_2}{v'_1 \\ \vdots \\ v'_p} \\
                    m_3 = \mathrm{Extend}(m_2, ((l_1, v_1), …, (l_n, v_n))) \\
                    \mm{m_3}{i}{m_4}{\iReturn{v}} \\
                    m_5 = \mathrm{Pop}(m_4) \\
                    m_6 = \mathrm{Cleanup}(m_5) \\
                  }
                  {\mm{m_0}{f (v_1, …, v_n)}{m_6}{v}}
    }
    {Seq}{ \irule{Seq}
             { \mmi{m}{i}{m'}{\iPass} }
             { \mmi{m}{(i;i')}{m'}{i'} }
    }
    {Pass}{ \irule{Pass}
              { }
              { \mmi{m}{(\iPass;i)}{m}{i} }
    }
    {Exp}{ \irule{Exp}
             { }
             { \mmi{m}{v}{m}{\iPass} }
    }
    {If-False}{ \irule{If-False}
                  { }
                  { \mmi{m}{\iIf{0}{i_t}{i_f}}{m}{i_f} }
    }
    {If-True}{ \irule{If-True}
                 { v ≠ 0 }
                 { \mmi{m}{\iIf{v}{i_t}{i_f}}{m}{i_t} }
    }
    {While}{ \irule{While}
               { }
               { \mmi{m}{\iWhile{e}{i}}{m}{\iThen{e}{i;\iWhile{e}{i}}}
               }
    }
    {Return}{ \irule{Return}
              { }
              { \mmi{m}{\iReturn{v};i}{m}{\iReturn{e}} }
    }
    {Ph-Struct}{ \irule{Ph-Struct}
                   { }
                   {\ph{m}{\tstruct{s} \{ … \}}{m}}
    }
    {Ph-Exp}{ \irule{Ph-Exp}
                {\mm{m}{e}{m'}{v}}
                {\ph{m}{e}{m'}}
            }
    {Ph-Var}{ \irule{Ph-Var}
                {\mm{m}{e}{m'}{v}}
                {\ph{(s,g)}{x = e}{(s, (x, v)::g)}}
    }
    {Ph*-Nil}{ \irule{Ph*-Nil}
                 { }
                 {\phstar{m}{[]}{m}}
    }
    {Ph*-Cons}{ \irule{Ph*-Cons}
                  {\ph{m}{p}{m'} \\
                   \phstar{m'}{ps}{m''}}
                  {\phstar{m}{p::ps}{m''}}
    }
    {END}
}

\def\figctx{
  \gramlr{Contextes}{
  \begin{align*}
  C \gramisa & C_L
  \\ \gramor & \ctxOp{C}{e}
  \\ \gramor & \ctxOp{v}{C}
  \\ \gramor & \ctxUnOp{C}
  \\ \gramor & \ctxSet{C}{e}
  \\ \gramor & \ctxSet{φ}{C}
  \\ \gramor & \{ l_1:v_1 ; … ; l_i:C ; … ; l_n:e_n \}
  \\ \gramor & [ v_1 ; … ; C ; … ; e_n ]
  \\ \gramor & C (e_1, …, e_n)
  \\ \gramor & f (v_1, …, C, …, e_n)
  \\
  \\
  C_L \gramisa & \ctxEmpty
  \\   \gramor & \ctxLvDeref{C_L}
  \\   \gramor & \ctxLvField{C_L}{l}
  \\   \gramor & \ctxLvIndex{C_L}{e}
  \\   \gramor & \ctxLvIndex{φ}{C}
  \\
  \\
  C_I \gramisa & C_I;i
  \\   \gramor & \iIf{C}{i_1}{i_2}
  \\   \gramor & \iReturn{C}
  \\   \gramor & C
  \end{align*}
  }
}


\usetheme{SAS}
\usetikzlibrary{arrows}
\usetikzlibrary{calc}
\usetikzlibrary{patterns}
\usetikzlibrary{shapes}
\usepackage{../gramster}
\usepackage{../pelagiccod}
\usepackage{booktabs}

\usepackage{amsmath,amsthm, amssymb, latexsym}
\usefonttheme[onlymath]{serif}
%\boldmath

\def\colK#1{\colorbox{MellonPink}{\strut $ #1 $}}
\def\colU#1{\colorbox{MellonGreen}{\strut $ #1 $}}
\colorlet{codeoutcol}{MellonYellow!60}
\def\codeout#1{
    \begin{center}
    \fboxsep=2mm
    \colorbox{codeoutcol}{\BUseVerbatim{#1}}
    \end{center}
}

\usepackage{fancyvrb}

\title{Analyse de sécurité de logiciels système par typage statique}
\subtitle{Application au noyau Linux}
\author{Étienne Millon}
\date{10 juillet 2014}

\begin{document}

\AtBeginSubsection[]
{
  \addtocounter{framenumber}{-1}
  \begin{frame}<beamer>
   \tableofcontents[currentsubsection]
 \end{frame}
}

\frame{\titlepage}

\frame{\tableofcontents}

\section{Le problème des pointeurs}

\subsection{L'arnaque du serrurier}

\begin{frame}

\begin{itemize}
\item
  j'appelle un serrurier
\item
  il peut m'ouvrir la porte
\item
  et si je lui dis que j'habite chez mon voisin? à la banque?
\end{itemize}

\end{frame}
\begin{frame}

\begin{itemize}
\item problème de confiance
\item il ne doit pas appliquer ses droits…
\item …mais ceux de l'appelant
\item il doit vérifier que j'ai bien les droits sur la maison
\end{itemize}

\end{frame}

\subsection{Isolation entre utilisateur et noyau}

\begin{frame}{Espaces utilisateur et noyau}
    \begin{itemize}
        \item espace utilisateur: programmes (navigateur, …)
        \item espace noyau: pilotes
        \item matériel
    \end{itemize}
\end{frame}

% Memory zone
%
% #1 - start
% #2 - end
% #3 - color
\newcommand{\mzone}[3]{
  \path[#3] (#1,0) rectangle (#2,1);
}

% Address label
%
% #1 - x position
% #2 - text
\newcommand{\alabel}[2]{
  \path (#1,1) -- ++(0,0.3) node [pos=1] {\small \tt #2};
}
\newcommand{\addrlabel}[2]{
  \draw[<-] (#1,-0.1) -- ++(0,-0.3) node [auto,pos=1] {\small \tt #2};
  \draw[pattern=north east lines] (#1,0) rectangle ++(0.2,1);
}
\long\def\memzones{

  % exec
  \mzone{0.5}{1}{user}

  % lib
  \mzone{2.5}{3.2}{user}

  % stack
  \mzone{3.7}{4}{user}

  % stack
  \mzone{5}{5.5}{user}

  % kernel
  \mzone{6}{8}{kernel}

  % contour
  \draw (0,0) rectangle (8,1);

  \alabel{0}{0}
  \alabel{6}{3 Go}
  \alabel{8}{4 Go}

}

\begin{frame}[fragile]{Séparation}
\begin{tikzpicture}
  [user/.style={fill=MellonGreen}
  ,kernel/.style={fill=MellonPink}
  ]

  \memzones{}

  \node at (2, -0.7) {Programme};
  \node at (7, -0.7) {Noyau};

  \draw[>=triangle 45,<->,ultra thick] (3.8, -0.7) -- node[auto] {Appels système} ++(2, 0);
  \draw[>=triangle 45,<->,ultra thick] (7, -1) -- node[auto,swap] {Instructions
  privilégiées} ++(0, -1) node (Htop) {};

  \path (Htop) -- ++(-1,-1) node (Htopleft) {};
  \draw (Htopleft) rectangle ++(2, -2);

  \foreach \x in {0,...,9} {
      \draw ($ (Htopleft) + (0.1 + 0.2*\x,  0) $) -- ++(0,  0.3);
      \draw ($ (Htopleft) + (0.1 + 0.2*\x, -2) $) -- ++(0, -0.3);
      \draw ($ (Htopleft) + (0, -0.2*\x - 0.1) $) -- ++(-0.3, 0);
      \draw ($ (Htopleft) + (2, -0.2*\x - 0.1) $) -- ++( 0.3, 0);
   }

  \path (Htop) -- ++(0, -2) node { Matériel };

\end{tikzpicture}
\end{frame}

\begin{frame}
    \insertcode{gettimeofday.c}

\begin{tikzpicture}
  [user/.style={fill=MellonGreen}
  ,kernel/.style={fill=MellonPink}
  ]

  \memzones{}
  \addrlabel{2.7}{tv}
\end{tikzpicture}
\end{frame}

\begin{frame}
    \insertcode{gettimeofday-bad.c}

\begin{tikzpicture}
  [user/.style={fill=MellonGreen}
  ,kernel/.style={fill=MellonPink}
  ]

  \memzones{}
  \addrlabel{6.7}{tv}
\end{tikzpicture}
\end{frame}

\begin{frame}
\begin{itemize}
    \item il faut vérifier tous les pointeurs
    \item fonctions \texttt{copy\_from\_user}, \texttt{copy\_to\_user}
    \item test + \texttt{memcpy}
    \item ou retournent une erreur ($< 0$)
    \item but: détecter les cas où on devrait utiliser \texttt{copy\_*\_user}
\end{itemize}
\end{frame}

\section{Formalisation}

\subsection{\langname}

\begin{frame}{Les problèmes de C}
    \begin{itemize}
        \item casts
        \item arithmétique de pointeurs
        \item peu de garanties
    \end{itemize}
\end{frame}

\begin{frame}
  \begin{align*}
  \gramdef{Expressions}{e}
                 { c               }{ Constante }
                 { \opun~e         }{ Opération unaire }
                 { e~\opbin~e      }{ Opération binaire }
                 { lv              }{ Accès mémoire }
                 { lv ← e          }{ Affectation }
                 { \& lv           }{ Pointeur }
                 { f               }{ Fonction }
                 { e (e_1, …, e_n) }{ Appel de fonction }
                 { \eStruct{
                    l_1 : e_1
                    ; …
                    ; l_n : e_n }  }{ Structure }
                 { \eArray{e_1 ;…; e_n} }{ Tableau }
                 {END}
  \end{align*}
\end{frame}

\begin{frame}
  \begin{align*}
  \gramdef{Instructions}{i}
                 { \iPass          }{Instruction vide}
                 { i;i             }{Séquence}
                 { e               }{Expression}
                 { \iDecl{x}{e}{i} }{Déclaration de variable}
                 { \iIf{e}{i}{i}   }{Alternative}
                 { \iWhile{e}{i}   }{Boucle}
                 { \iReturn{e}     }{Retour de fonction}
                 {END}
  \end{align*}
\end{frame}
\subsection{Sémantique}

\begin{frame}{Sémantique impérative}

\begin{itemize}
\item
  Système de transitions entre états
\item
  expression + mémoire
\item
  $\mm{m}{e}{m'}{e'}$
\item expressions $→$ valeurs
\end{itemize}

\end{frame}

\begin{frame}
    \begin{itemize}
        \item valeurs gauches (\emph{lvalues}): \texttt{t[20] = 1}
        \item valeurs gauches $→$ chemins $φ$
        \item formes:
            \begin{itemize}
                \item \texttt{x} $→$ \texttt{(n, x)}
                \item \texttt{lv.l} $→ φ.l$
                \item \texttt{lv[e]} $→ φ[n]$
                \item \texttt{*e} $→ φ$
            \end{itemize}
        \item $\mm{m}{lv ← e}{m'}{v}$ ?
        \item ex: $x.f[5] ← 3$
        \item Mise à jour complète : $m' = m[ x ← 0]$
        \item Mise à jour partielle : $m' = m[x.f ← 0]$
    \end{itemize}
\end{frame}

\begin{frame}{Structure de la mémoire}

\begin{itemize}
\item $m = (s, g) ∈ \sMem$
\item $\sMem = \sList{\sFrame} × \sFrame$
\item liste de cadres de pile, cadre de globales
\item cadre: $c = (x_1 ↦ v_1, …, x_n ↦ v_n) ∈ \sFrame$
\item Comment exprimer:
  \begin{itemize}
  \item $m[φ]$
  \item $m[φ←v]$
  \end{itemize}
\end{itemize}
\end{frame}

\begin{frame}{À la main}

\begin{align*}
(s, g)[\texttt{Local}~(2, x).f[5] ← 3] &= (s', g) \\
                   s'_i &= \begin{cases}
                             s_i & \mbox{ si } i ≠ 2\\
                             c'  & \mbox{ sinon} \\
                           \end{cases}\\
                         c &= s_{2} \mbox{ où } c_i &= (x_i, v_i) \\
                   c'_i &= (x_i, v'_i) \\
                   v'_i &= \begin{cases}
                              v_i &  \mbox{ si } i ≠ 1 \\
                              3   &  \mbox{ sinon }\\
                           \end{cases} \\
                    … & …
\end{align*}
\end{frame}

\subsection{Lentilles}

\begin{frame}{Lentilles}

\begin{itemize}
\item
  Relation entre objet et sous-objet
\item
  $ℒ ∈ \setLens{R}{A}$
\item
  $ℒ = \begin{cases}           \mathrm{get}_ℒ : R → A \\           \mathrm{put}_ℒ : (A × R) → R \\          \end{cases}$
\item
  $\lensGet{ℒ}{\lensNodeBig{green!30,ultra thick}{blue!30, ultra thick}} = \lensInner{blue!30, ultra thick}$
\item
  $\lensPut{ℒ}{\lensInner{red!30, ultra thick}}{\lensNodeBig{green!30, ultra thick}{blue!30, ultra thick}} = \lensNodeBig{green!30, ultra thick}{red!30, ultra thick}$
\end{itemize}

\end{frame}

\begin{frame}

\begin{center}
\definecolor{figcola}{HTML}{EFEF52}
\definecolor{figcolb}{HTML}{52EFA1}
\definecolor{figcolc}{HTML}{52A1EF}
\definecolor{figcold}{HTML}{EFA152}
\begin{tikzpicture}
[node distance=2.5cm
,bignode/.style={draw,shape=rectangle,minimum size=1.5cm,fill=figcola}
,smallnode/.style={draw,shape=circle,minimum size=1cm,fill=figcolb}
,trinode/.style={draw,shape=regular polygon,regular polygon sides=3}
,>=triangle 45
,ultra thick
]
\node[bignode] (A) {};
\node[right of=A,smallnode] (B) {};
\node[right of=B,trinode,fill=figcolc] (C) {};
\node[below of=B,smallnode] (D) {};
\node[below of=A,node distance=5cm,bignode] (E) {};

\draw[->] (A) to node[auto] {$\mathrm{get}_{ℒ_1}$} (B);
\draw[->] (B) to (D);
\draw[->] (A) to (E);
\draw (C) |- node[auto,near end,swap] {$\mathrm{put}_{ℒ_2}$} ($(B)!0.5!(D)$);
\draw (D) |- node[auto,near end,swap] {$\mathrm{put}_{ℒ_1}$} ($(A)!0.75!(E)$);

\node[smallnode] at (A) {};
\node[trinode,fill=figcold] at (A) {};

\node[trinode,fill=figcold] at (B) {};

\node[trinode,fill=figcolc] at (D) {};

\node[smallnode] at (E) {};
\node[trinode,fill=figcolc] at (E) {};

\node[left of=A, node distance=4cm,bignode] (GA) {};
\node[below of=GA,smallnode] (GB) {};
\node[below of=GB,trinode,fill=figcold] (GC) {};

\node[smallnode] at (GA) {};

\node[trinode, fill=figcold] at (GA) {};
\node[trinode, fill=figcold] at (GB) {};

\draw [->] (GA) to node[auto] {$\mathrm{get}_{ℒ_1}$} (GB);
\draw [->] (GB) to node[auto] {$\mathrm{get}_{ℒ_2}$} (GC);

\end{tikzpicture}

\vspace{-1cm}\hspace{7cm} \fbox{$ℒ_1 \ggg ℒ_2$}

\end{center}

\end{frame}

\begin{frame}{Lentilles en sémantique}
    \[Φ ∈ \sPath → \setLens{Mem}{Val}\]

    Pour $φ = (2, x).f[5]$:

    \begin{align*}
    Φ(φ) &= \texttt{fst} \\
         &\ggg I(2) \\
         &\ggg L(x) \\
         &\ggg F(f) \\
         &\ggg T(5)
    \end{align*}
\end{frame}

\begin{frame}{Lentilles en sémantique}

\begin{itemize}
\itemsep1pt\parskip0pt\parsep0pt
\item
  par tableau

  \begin{itemize}
  \itemsep1pt\parskip0pt\parsep0pt
  \item
    $T(n) ∈ \setLens{Val}{Val}$
  \item
    $T(3) \equiv [1, 2, 3, \circled{5}]$
  \end{itemize}
\item
  par champ

  \begin{itemize}
  \itemsep1pt\parskip0pt\parsep0pt
  \item
    $F(l) ∈ \setLens{Val}{Val}$
  \item
    $F(y) \equiv \{ x: 0; y: \circled{-3} \}$
  \end{itemize}
\end{itemize}

\end{frame}

\begin{frame}{Exemple : par variable ($V$)}

\[V(v) ∈ \setLens{Mem}{Val}\]

\begin{itemize}
\itemsep1pt\parskip0pt\parsep0pt
\item
  Pour $m ∈ Mem$, $m = ((s_1, …, s_n), g)$
\end{itemize}

\begin{align*}
V(\texttt{Global}~x) &= \texttt{snd} \ggg L(x) \\
V(\texttt{Local}~(n, x)) &= \texttt{fst} \ggg I(n) \ggg L(x) \\
\end{align*}

\end{frame}

\begin{frame}{Exemple : par chemin ($Φ$)}

\[Φ(φ) ∈ \setLens{Mem}{Val}\]

\begin{align*}
    Φ(v)    &= V(v) \\
  Φ(φ.l)    &= Φ(φ) \ggg F(l) \\
  Φ(φ[i])   &= Φ(φ) \ggg T(i) \\
  \end{align*}

\end{frame}

\subsection{Types pointeurs}

\begin{frame}{Deux types de pointeurs}

\begin{itemize}
\item séparer deux types de pointeurs
\item pointeurs noyau : $\ptrK{t}$

  \begin{itemize}
  \item valeur fixée à la compilation ou par le runtime
  \item ex: \texttt{\&x}, \texttt{malloc(n)}
  \item peuvent être déréférencés (\texttt{*p}, \texttt{memcpy()})
  \end{itemize}
\item pointeurs utilisateur : $\ptrU{t}$

  \begin{itemize}
  \item valeur provient d'un utilisateur non privilégié
  \item ex: paramètres d'appels systèmes
  \item doivent être vérifiés dynamiquement!
  \end{itemize}
\end{itemize}

\end{frame}

\begin{frame}{Règles de typage}
    \begin{mathpar}
       \irule{Lv-Deref}
         { Γ ⊢ \colK{e} : \colK{\ptrK{t}} }
         { Γ ⊢ \colK{*~e} : t }

       \irule{Addr}
         { Γ ⊢ lv : t }
         { Γ ⊢ \& lv : \colK{\ptrK{t}} }

       \irule{User-Get}
         { Γ ⊢ \colK{e_d} : \colK{\ptrK{t}}
        \\ Γ ⊢ \colU{e_s} : \colU{\ptrU{t}}
         }
         { Γ ⊢ \uGet{\colK{e_d}}{\colU{e_s}} : \tInt }

       \irule{User-Put}
         { Γ ⊢ \colU{e_d} : \colU{\ptrU{t}}
        \\ Γ ⊢ \colK{e_s} : \colK{\ptrK{t}}
         }
         { Γ ⊢ \uPut{\colU{e_d}}{\colK{e_s}} : \tInt }
    \end{mathpar}
\end{frame}

\begin{frame}{Sûreté du typage}

\begin{theorem}[Progrès]
  Supposons que $Γ ⊢ e : t$. Soit $m$ un état mémoire tel que $\mcomp{Γ}{m}$.
  Alors l'un des cas suivants est vrai:

\begin{itemize}
  \item $∃ v ≠ Ω, e = v$
  \item $∃ (e', m'), \mm{m}{e}{m'}{e'}$
  \item $∃ Ω ∈ \{\serr{div},\serr{array},\serr{ptr}\}, \msi{m}{e} → Ω$
\end{itemize}
\end{theorem}
\end{frame}

\begin{frame}{Sûreté du typage}

\begin{theorem}[Préservation]
    Si $Γ ⊢ e : t$ et $\mm{m}{e}{m'}{v}$,
    alors $\mcomp{Γ}{\cCleanup{m'}}$ et $m' ⊧ v : τ$ où $\tComp{τ}{t}$.
\end{theorem}
\end{frame}

\section{Expérimentation}

\subsection{Analyseur statique}

\begin{frame}

\begin{itemize}
\item 1600 lignes d'OCaml
\item utilise \newspeak (langage + compilateur)
\end{itemize}

\end{frame}

\begin{frame}{Mode de fonctionnement}

  \begin{itemize}
  \item compile vers un \texttt{.npk}
  \item enlève les types C
  \item infère des types plus sémantiques
  \end{itemize}
\end{frame}

\begin{frame}{Traduction: C} \insertcode{tc-c.c} \end{frame}
\begin{frame}{Traduction: sans les types} \insertcode{tc-ml.c} \end{frame}
\begin{frame}{Traduction: types inférés} \insertcode{tc-ty.c} \end{frame}

\begin{frame}
    \centering
      \begin{tabular}{ll}
          \toprule
          C/\newspeak & \langname \\
          \midrule
          \texttt{int} & $\tInt$ \\
          \texttt{unsigned long} & $\tInt$ \\
          \texttt{int *} (utilisé \texttt{x[i]}) & $\tInt~[~]$ \\
          \texttt{int *} (utilisé \texttt{*x}) & $\ptrK{\tInt}$ \\
          \texttt{int *} (argument syscall) & $\ptrU{\tInt}$ \\
          \bottomrule
      \end{tabular}
\end{frame}

\subsection{Utilisation sur le noyau Linux}

\begin{frame}{Analyse de 2 bugs}

    \begin{itemize}
        \item Radeon KMS
        \item \texttt{ptrace} Blackfin
    \end{itemize}
\end{frame}

\begin{frame}{Code avec bug}
    % left bottom right top
    \includegraphics[trim=2300 1430 100 500,clip,width=\textwidth]{poster.pdf}
\end{frame}

\begin{frame}[fragile]{Résultat de l'algorithme d'inférence}
\begin{SaveVerbatim}{drmko}
05-drm.c:17#10 - Type clash between :
  KPtr (_a15)
  UPtr (_a8)
\end{SaveVerbatim}

\textbf{Erreur}:

\codeout{drmko}
\end{frame}

\begin{frame}{Code corrigé}
    % left bottom right top
    \includegraphics[trim=2300 990 100 1220,clip,width=\textwidth]{poster.pdf}
\end{frame}

\begin{frame}[fragile]{Résultat de l'algorithme d'inférence}

\textbf{Programme complètement annoté}:

\begin{SaveVerbatim}{drmok}
(06-drm-ok.c:17#6)^{
  Int tmp_cir!0;
  (06-drm-ok.c:17#6)^tmp_cir!0 <-
    copy_from_user
      ( &(value) : KPtr (d),
        value_ptr_UPtr (d) : UPtr (d),
        4 : Int); }
\end{SaveVerbatim}

\codeout{drmok}
\end{frame}

\section{Perspectives}

\subsection{Contributions}

\begin{frame}{Contributions}
\begin{itemize}
\item Un langage impératif bien typé
\item Une sémantique basée sur les lentilles
\item Un système de types abstraits
\item Un prototype d'analyseur statique
\end{itemize}
\end{frame}

\subsection{Analyse à plus grande échelle}

\begin{frame}
    \begin{itemize}
        \item Comment détecter les appels système?
        \item \texttt{SYSCALL\_DEFINEx}, …
        \item le reste: \texttt{ioctl}s, FS, …
    \end{itemize}
\end{frame}

\subsection{Types abstraits}

\begin{frame}{Types abstraits : exemples}
    \begin{itemize}
        \item extension: comment le type est représenté
        \item intention: à quoi il sert
        \item \texttt{int}, \texttt{int32\_t}
        \item \texttt{off\_t}, \texttt{size\_t}, \texttt{intptr\_t}
        \item Types "numéro de ligne", "numéro de colonne"
    \end{itemize}
\end{frame}

\subsection{Conclusion}

\begin{frame}
    Merci

    Questions?
\end{frame}

\end{document}
